\section{Monte Carlo}
\label{sec:simulations}
\begin{align*}
R^e_{i,t+1}&=\alpha(Z_{i,t})+\beta(Z_{i,t})^\top F_{t+1}+\varepsilon_{i,t+1},\\
Z_{i,t,k}&=2\Phi(S_{i,t,k})-1,\qquad
S_{i,t,k}=\rho S_{i,t-1,k}+\sqrt{1-\rho^2}\,\eta_{i,t,k},\quad k=1,2,3,\\
\alpha(z)&=A_\alpha\sqrt2\cos\left(\pi\sum_{k=1}^3z_k\right),\\
\beta_k(z)&=A_\beta\sum_{m=2}^\infty\frac{\cos(\pi m z_k)}{m^{7/2}\log(m+1)},\quad k=1,2,3.
\end{align*}
We assume $|\rho|<1$, $A_\alpha,A_\beta,\sigma_\varepsilon>0$, $\lambda^\top\lambda<1$,
$S_{i,0,k}\sim\mathcal N(0,1)$ and $\eta_{i,t,k}\overset{\mathrm{i.i.d.}}\sim\mathcal N(0,1)$,
independently across stocks, characteristics and dates. Thus $Z_{i,t}$ has independent
$\operatorname{Unif}(-1,1)$ coordinates, and the characteristic vectors are independent
across stocks at each date. Define
\begin{align*}
\alpha_t&=(\alpha(Z_{1,t}),\ldots,\alpha(Z_{N,t}))^\top,
&B_t&=(\beta(Z_{1,t}),\ldots,\beta(Z_{N,t}))^\top,\\
\mu_t&=\alpha_t+B_t\lambda,
&\Omega_F&=I_3-\lambda\lambda^\top,\\
\Sigma_t^0&=\mu_t\mu_t^\top+B_t\Omega_FB_t^\top+\sigma_\varepsilon^2I_N,
&D_t&=\kappa_0 I_N-\Sigma_t^0.
\end{align*}
Choose a constant $\kappa_0>\sup_{Z_t}\lambda_{\max}(\Sigma_t^0)$ and set
\[
h_t=\begin{cases}
D_t\mu_t/\sqrt{\mu_t^\top D_t\mu_t},&\mu_t\neq0,\\
0,&\mu_t=0.
\end{cases}
\]
Conditional on the decision-date information $\mathscr F_t$, the \emph{Gaussian generators}
are independent and satisfy
\begin{align*}
F^G_{t+1}\mid\mathscr F_t&\sim\mathcal N_3(\lambda,\Omega_F),\\
\varepsilon^G_{t+1}\mid\mathscr F_t&\sim
\mathcal N_N(0,\sigma_\varepsilon^2I_N+h_th_t^\top).
\end{align*}
Equivalently, with independent standard Gaussian vectors and scalars,
\begin{align*}
F^G_{t+1}&=\lambda+\Omega_F^{1/2}G^F_{t+1},\\
\varepsilon^G_{t+1}&=\sigma_\varepsilon G^\varepsilon_{t+1}+h_tG^0_{t+1},\\
G^F_{t+1}&\sim\mathcal N_3(0,I_3),\quad
G^\varepsilon_{t+1}\sim\mathcal N_N(0,I_N),\quad
G^0_{t+1}\sim\mathcal N(0,1).
\end{align*}
All Gaussian generators are independent across dates and of the characteristic process.
For bounded returns, replace each primitive $G\sim\mathcal N(0,1)$ in the
\emph{factor and return innovations} by the standardized truncated Gaussian
\[
\xi=\frac{G\mid\{|G|\leq a\}}{\sqrt{v_a}},\qquad
v_a=1-\frac{2a\varphi(a)}{2\Phi(a)-1},\qquad a=4.47.
\]
Use the resulting innovations in the preceding linear formulas to define the actual
$F_{t+1}$ and $\varepsilon_{t+1}$. They retain exactly the displayed conditional
means and covariance matrices, although after truncation they are no longer
literally Gaussian. Each return uses at most five independent primitive innovations
(three factors, one individual shock and one common shock), hence
\[
|R^e_{i,t+1}-\mu_{i,t}|
\leq\sqrt5\,\frac{a}{\sqrt{v_a}}
\sqrt{\operatorname{Var}(R^e_{i,t+1}\mid\mathscr F_t)}
<10\sqrt{\operatorname{Var}(R^e_{i,t+1}\mid\mathscr F_t)}.
\]
This enforces a ten-conditional-standard-deviation bound without clipping the
finished returns and thereby changing their conditional moments. Here
$\mathscr F_t$ contains the characteristics at $t$ and the history of
characteristics and realized innovations through $t$.
The resulting conditional return law has
\begin{align*}
\E[R^e_{t+1}\mid\mathscr F_t]&=\mu_t,\\
\Sigma_t:=\E[R^e_{t+1}(R^e_{t+1})^\top\mid\mathscr F_t]
&=\Sigma_t^0+h_th_t^\top.
\end{align*}
The characteristic distribution implies
\begin{align*}
\E[\alpha(Z)]&=0,&\E[\beta(Z)]&=0,&\E[\alpha(Z)\beta(Z)]&=0,\\
\E[\alpha(Z)^2]&=v_\alpha=A_\alpha^2,
&\E[\beta(Z)\beta(Z)^\top]&=v_\beta I_3,\\
v_\beta&=\frac{A_\beta^2}{2}\sum_{m=2}^\infty
\frac1{m^7\log^2(m+1)}.
\end{align*}
We use the Mat\'ern-$3/2$ kernel with a fixed lengthscale $\ell>0$,
\[
K(z,z')=\left(1+\frac{\sqrt3\|z-z'\|}{\ell}\right)
\exp\left(-\frac{\sqrt3\|z-z'\|}{\ell}\right),
\qquad\Hh_K\simeq H^3((-1,1)^3),
\]
where $\simeq$ denotes equality of function sets with equivalent norms.
As in the population problem, $R^p_{t+1}(W)=N^{-1}\sum_iW(Z_{i,t})R^e_{i,t+1}$ and
$Q(W)=\E[(1-R^p_{t+1}(W))^2]$.
\begin{proposition}[Population-optimal portfolio policy]
\label{prop:mc_optimal}
Under the specified DGP, the unique population-optimal scalar policy is
\[
W^\star(z)=\frac{N}{\kappa_0}
\left[\alpha(z)+\beta(z)^\top\lambda\right].
\]
It minimizes $Q(W)$ over $\Hh_K$ and also minimizes the conditional quadratic
loss at every date. The actual portfolio weight of stock $i$ is
$W^\star(Z_{i,t})/N$.
\end{proposition}
\begin{proof}
By construction, $D_t\succ0$. For $\mu_t\neq0$,
\begin{align*}
h_th_t^\top\mu_t
&=\frac{D_t\mu_t\mu_t^\top D_t\mu_t}
{\mu_t^\top D_t\mu_t}
=D_t\mu_t.
\end{align*}
The same equality is trivial if $\mu_t=0$. Consequently,
\[
\Sigma_t\mu_t=(\Sigma_t^0+h_th_t^\top)\mu_t
=\Sigma_t^0\mu_t+D_t\mu_t=\kappa_0\mu_t.
\]
For any scalar policy $W$, write
\[
w_t(W)=\frac1N(W(Z_{1,t}),\ldots,W(Z_{N,t}))^\top.
\]
The conditional loss is
\[
Q_t(W)=1-2w_t(W)^\top\mu_t+w_t(W)^\top\Sigma_tw_t(W).
\]
The proposed policy gives $w_t(W^\star)=\mu_t/\kappa_0$, and thus
$\Sigma_tw_t(W^\star)=\mu_t$. Completing the square,
\begin{align*}
Q_t(W)-Q_t(W^\star)
&=[w_t(W)-w_t(W^\star)]^\top\Sigma_t
[w_t(W)-w_t(W^\star)]\\
&\geq\frac{\sigma_\varepsilon^2}{N^2}
\sum_{i=1}^N[W(Z_{i,t})-W^\star(Z_{i,t})]^2.
\end{align*}
Taking expectations and using the common stock-characteristic distribution,
\[
Q(W)-Q(W^\star)\geq\frac{\sigma_\varepsilon^2}{N}
\E[(W(Z)-W^\star(Z))^2].
\]
By Proposition~\ref{prop:mc_sobolev}, $W^\star\in\Hh_K$, so the lower bound is
attained by an admissible scalar policy. It is strictly positive for every
policy differing from $W^\star$ on a set of positive $P_Z$-measure, establishing
existence and uniqueness (up to $P_Z$-almost sure equality).
\end{proof}
\begin{proposition}[Optimal Sharpe ratio]
\label{prop:mc_sharpe}
Set $L=\lambda^\top\lambda$ and
\[
p^\star=\frac{N}{\kappa_0}(v_\alpha+v_\beta L).
\]
Then $0<p^\star<1$ and
\begin{align*}
\E[R^p_{t+1}(W^\star)]&=\E[(R^p_{t+1}(W^\star))^2]=p^\star,\\
Q(W^\star)&=1-p^\star,\\
\operatorname{Var}(R^p_{t+1}(W^\star))&=p^\star(1-p^\star),\\
\SR(W^\star)&=\sqrt{\frac{p^\star}{1-p^\star}},\\
\SR_{\mathrm{ann,marg}}(W^\star)&=\sqrt{12}\sqrt{\frac{p^\star}{1-p^\star}}.
\end{align*}
This is the maximum population Sharpe ratio over admissible scalar policies.
\end{proposition}
\begin{proof}
The stock weights satisfy $w_t^\star=\mu_t/\kappa_0$ and
$\Sigma_t\mu_t=\kappa_0\mu_t$. Therefore
\begin{align*}
\E[R^p_{t+1}(W^\star)\mid\mathscr F_t]
&=(w_t^\star)^\top\mu_t=\frac{\|\mu_t\|^2}{\kappa_0},\\
\E[(R^p_{t+1}(W^\star))^2\mid\mathscr F_t]
&=(w_t^\star)^\top\Sigma_tw_t^\star=\frac{\|\mu_t\|^2}{\kappa_0}.
\end{align*}
Cross-sectional linearity, orthogonality and the common marginal distribution give
\begin{align*}
\E[\|\mu_t\|^2]
&=\sum_i\E[(\alpha(Z_{i,t})+\beta(Z_{i,t})^\top\lambda)^2]\\
&=N(v_\alpha+v_\beta L).
\end{align*}
Thus both unconditional moments equal $p^\star$. The mean is positive because
$A_\alpha>0$. The conditional covariance matrix of returns satisfies
$\Sigma_t-\mu_t\mu_t^\top\succeq\sigma_\varepsilon^2I_N\succ0$.
If $\mu_t\neq0$, evaluating its quadratic form at $\mu_t$ gives
\[
0<\kappa_0\|\mu_t\|^2-\|\mu_t\|^4,
\qquad\frac{\|\mu_t\|^2}{\kappa_0}<1.
\]
Consequently $0<p^\star<1$. Expanding $Q(W^\star)$ and subtracting the
square of the mean from the second moment yields the stated loss, variance
and Sharpe formulas.
For any admissible nonzero payoff $Y=R^p_{t+1}(W)$, scalar rescaling gives
\[
\inf_{c\in\mathbb R}\E[(1-cY)^2]
=1-\frac{\E[Y]^2}{\E[Y^2]}
=\frac1{1+\SR(Y)^2}.
\]
Since the left side is at least $Q(W^\star)$ and $W^\star$ has a positive mean,
its Sharpe is maximal. Multiplication by $\sqrt{12}$ is the marginal monthly
to annual convention; it does not assert a serial-independence formula for
compounded annual returns.
\end{proof}
\begin{proposition}[Calibration of the DGP]
\label{prop:mc_calibration}
Define
\[
\kappa_\beta=\frac12\sum_{m=2}^\infty\frac1{m^7\log^2(m+1)},
\qquad v_\beta=\kappa_\beta A_\beta^2.
\]
For a desired optimal expected payoff $p_0\in(0,1)$ and fixed
$\kappa_0,A_\beta,\lambda$, the alpha amplitude is
\[
A_\alpha=\left[\frac{\kappa_0p_0}{N}
-\kappa_\beta A_\beta^2\lambda^\top\lambda\right]^{1/2}.
\]
Alternatively, for fixed $A_\alpha,\kappa_0$ and $L=\lambda^\top\lambda>0$,
\[
A_\beta=\left[
\frac{\kappa_0p_0/N-A_\alpha^2}{\kappa_\beta L}
\right]^{1/2}.
\]
Equivalently, for fixed amplitudes,
\[
\kappa_0=\frac{N}{p_0}
\left(A_\alpha^2+\kappa_\beta A_\beta^2L\right).
\]
These calibrations require positive radicands and
$\kappa_0>\sup_{Z_t}\lambda_{\max}(\Sigma_t^0)$.
For $\SR_{\mathrm{ann,marg}}^\star=3$, the equivalent restriction is
\[
p_0=\frac37,\qquad
\kappa_0=\frac{7N}{3}
\left(A_\alpha^2+\kappa_\beta A_\beta^2L\right).
\]
\end{proposition}
\begin{proof}
Proposition~\ref{prop:mc_sharpe} implies
\[
p_0=\frac{N}{\kappa_0}
\left(A_\alpha^2+\kappa_\beta A_\beta^2L\right).
\]
Solving this identity for $A_\alpha$, $A_\beta$ or $\kappa_0$ gives the
three expressions. Moreover,
\[
3=\sqrt{12}\sqrt{\frac{p_0}{1-p_0}}
\quad\Longleftrightarrow\quad p_0=\frac37.
\]
The admissible set is nonempty. To see this, hold $A_\alpha>0$, $N$ and
$0<L<1$ fixed. As $A_\beta,\sigma_\varepsilon\downarrow0$,
\[
\kappa_0\longrightarrow\frac{7N}{3}A_\alpha^2,
\qquad
\sup_{Z_t}\lambda_{\max}(\Sigma_t^0)
\leq2NA_\alpha^2+o(1).
\]
Because $7/3>2$, sufficiently small strictly positive $A_\beta$ and
$\sigma_\varepsilon$ satisfy the strict covariance condition while
retaining a nonzero factor component and a strictly positive idiosyncratic
variance.
\end{proof}
\begin{proposition}[Sobolev regularity of the optimal policy]
\label{prop:mc_sobolev}
Let $\mathcal Z=(-1,1)^3$. If $A_\beta>0$ and
$\lambda^\top\lambda>0$, then
\[
W^\star\in H^3(\mathcal Z)=\Hh_K,
\qquad
W^\star\notin H^{3+\delta}(\mathcal Z)
\quad\text{for every }\delta>0.
\]
\end{proposition}
\begin{proof}
The function $\alpha$ is a trigonometric polynomial, hence belongs to
$H^r(\mathcal Z)$ for every finite $r\geq0$. Each $\beta_k$ is a two-periodic
cosine series with nonzero Fourier coefficients
\[
\widehat\beta_k(\pm me_k)
=\frac{A_\beta}{2m^{7/2}\log(m+1)},\qquad m\geq2.
\]
Consequently, for any $r\geq0$,
\[
\|\beta_k\|_{H^r(\mathbb T^3)}^2
\asymp A_\beta^2\sum_{m=2}^\infty
\frac{m^{2r-7}}{\log^2(m+1)}.
\]
At $r=3$ this is a convergent series
$\sum_{m\geq2}[m\log^2(m+1)]^{-1}$, whereas for every
$\delta>0$ the corresponding series at $r=3+\delta$ is
$\sum_{m\geq2}m^{2\delta-1}/\log^2(m+1)=\infty$.
Thus $\beta_k\in H^3(\mathbb T^3)$ but
$\beta_k\notin H^{3+\delta}(\mathbb T^3)$ for every $\delta>0$.
The policy in Proposition~\ref{prop:mc_optimal} is a finite linear
combination of $\alpha$ and the $\beta_k$, and therefore belongs to $H^3$.
Choose $k$ such that $\lambda_k\neq0$. The Fourier support of $\alpha$
is $\{\pm(1,1,1)\}$, whereas the support of $\beta_k$ is
$\{\pm me_k:m\geq2\}$; the supports of different loadings are also
disjoint. Hence
\[
\widehat W^\star(\pm me_k)
=\frac{NA_\beta\lambda_k}{2\kappa_0m^{7/2}\log(m+1)},
\]
and for every $\delta>0$,
\[
\|W^\star\|_{H^{3+\delta}(\mathbb T^3)}^2
\geq c\sum_{m=2}^\infty
\frac{m^{2\delta-1}}{\log^2(m+1)}=\infty.
\]
The restriction to the cube has the same sharp threshold. Indeed, the
periodic series and its derivatives of orders at most two are continuous,
and their traces match across opposite cube faces. If the restriction were
in $H^{3+\delta}(\mathcal Z)$, choose
$0<\delta_0<\min\{\delta,1/2\}$. The Sobolev gluing criterion would put
its periodic extension in $H^{3+\delta_0}(\mathbb T^3)$, contradicting
the preceding Fourier divergence.
Finally, on the bounded Lipschitz domain $\mathcal Z\subset\mathbb R^3$,
the Mat\'ern-$3/2$ RKHS is precisely $H^{3/2+3/2}(\mathcal Z)=H^3(\mathcal Z)$
as a function set, with equivalent norms. This proves both claims.
\end{proof}
\begin{proposition}[Verification of the economy assumptions]
\label{prop:mc_assumptions}
Fix $N$, $A_\alpha,A_\beta,\sigma_\varepsilon>0$, a nonzero $\lambda$
with $\lambda^\top\lambda<1$, a fixed Mat\'ern-$3/2$ kernel and
$\kappa_0>\sup_{Z_t}\lambda_{\max}(\Sigma_t^0)$.
Let $\rho$ vary in a compact interval $[-\bar\rho,\bar\rho]$,
$\bar\rho<1$. The resulting family of economies satisfies all six
conditions (E1)--(E6) in the portfolio-learning theorem with common
constants, $b=2$ and any prescribed $\gamma>2$. More precisely,
\begin{align*}
\beta_\varepsilon(n)&\leq C_\gamma n^{-\gamma},\\
\|X_t\|_{\Hh_K}&\leq B_X,\qquad
\|W^\star\|_{\Hh_K}\leq R,\\
\E[(1-R^p_{t+1}(W^\star))R^p_{t+1}(V)\mid\mathscr F_t]&=0
\quad\text{for every }V\in\Hh_K,\\
c_\mu j^{-2}&\leq\mu_j(\Sigma_{\Hh_K})\leq C_\mu j^{-2},\\
0<\underline S&\leq\SR(W^\star)\leq\overline S<\infty.
\end{align*}
The class also satisfies the order-$r=1$ source condition and the bounded
conditional Bernstein requirement for the kernel-weighted regression residual.
\end{proposition}
\begin{proof}
\emph{(E1): stationarity and mixing.}
The joint latent state $S_t=(S_{i,t,k})_{i,k}$ is a finite-dimensional
stationary Gaussian AR(1) Markov chain. Its $n$-step transition law is
$\mathcal N(\rho^n x,(1-\rho^{2n})I_{3N})$. Using Gaussian relative
entropy and the total-variation bound gives
\[
\beta_S(n)\leq
\left[-\frac{3N}{4}\log(1-\rho^{2n})\right]^{1/2}
\leq C\bar\rho^n.
\]
The primitive factor and idiosyncratic generators are independent
of the state process and of their own past. The financial observation
$(Z_t,R^e_{t+1})$ is a measurable function of the current state and
these fresh innovations; it is therefore stationary and geometrically
absolutely regular. Since
$\sup_{n\geq1}n^\gamma\bar\rho^n<\infty$ for every fixed
$\gamma>0$, the required polynomial inequality holds for any
$\gamma>2=b$. The bound is uniform over the stated interval for $\rho$.
\par\noindent\emph{(E2): bounded managed features.}
Let $\mathcal B=\sum_{m=2}^\infty[m^{7/2}\log(m+1)]^{-1}$.
Then, uniformly over $z$,
\[
|\alpha(z)|\leq\sqrt2 A_\alpha,
\qquad\|\beta(z)\|\leq\sqrt3 A_\beta\mathcal B.
\]
The matrices $D_t$ are positive definite and bounded by $\kappa_0I_N$.
For $\mu_t\neq0$, the inequality $D_t^2\preceq\|D_t\|D_t$ yields
\[
\|h_t\|^2
=\frac{\mu_t^\top D_t^2\mu_t}{\mu_t^\top D_t\mu_t}
\leq\|D_t\|\leq\kappa_0.
\]
The same bound holds when $\mu_t=0$. The conditional standard deviation
of return $i$ is therefore bounded by
\[
\sigma_{i,t}:=
\sqrt{\beta(Z_{i,t})^\top\Omega_F\beta(Z_{i,t})
+\sigma_\varepsilon^2+h_{i,t}^2}
\leq\sqrt{3A_\beta^2\mathcal B^2+\sigma_\varepsilon^2+\kappa_0}.
\]
The ten-sigma bound and the bound on $\mu_{i,t}$ imply
\begin{align*}
|R^e_{i,t+1}|
&\leq\sqrt2 A_\alpha+\sqrt3 A_\beta\mathcal B\|\lambda\|\\
&\quad+10\sqrt{3A_\beta^2\mathcal B^2+\sigma_\varepsilon^2+\kappa_0}
=:B_R.
\end{align*}
For $X_t=N^{-1}\sum_iR^e_{i,t+1}K(\cdot,Z_{i,t})$ and
$K(z,z)=1$, the reproducing property gives
\[
\|X_t\|_{\Hh_K}\leq\frac1N\sum_i
|R^e_{i,t+1}|\sqrt{K(Z_{i,t},Z_{i,t})}
\leq B_R=:B_X.
\]
Thus (E2) holds with $\|X_t\|_{\Hh_K}^2\leq B_X^2$.
\par\noindent\emph{(E3): population attainment, norm and source condition.}
Propositions~\ref{prop:mc_optimal} and~\ref{prop:mc_sobolev}
show that the minimum is attained and $W^\star\in\Hh_K$. Furthermore,
\[
\|W^\star\|_{\Hh_K}
\leq\frac{N}{\kappa_0}
\left(\|\alpha\|_{\Hh_K}
+\|\lambda\|\Bigl[\sum_{k=1}^3\|\beta_k\|_{\Hh_K}^2\Bigr]^{1/2}\right)
=:R<\infty.
\]
This bound does not depend on $\rho$. Strict convexity from
Proposition~\ref{prop:mc_optimal} also makes this the minimum-norm
population minimizer. Membership in $\Hh_K$ verifies the source
condition of order $r=1$; it does not imply a higher source order.
\par\noindent\emph{(E4): conditional pricing.}
For every $V\in\Hh_K$ and $v_t=(V(Z_{1,t}),\ldots,V(Z_{N,t}))^\top/N$,
\begin{align*}
&\E[(1-R^p_{t+1}(W^\star))R^p_{t+1}(V)\mid\mathscr F_t]\\
&=v_t^\top\bigl(\mu_t-\Sigma_tw_t(W^\star)\bigr)\\
&=v_t^\top\left(\mu_t-\Sigma_t\frac{\mu_t}{\kappa_0}\right)=0.
\end{align*}
This proves the requested condition against the entire
$\mathscr F_t$, not merely the unconditional portfolio first-order condition.
The pricing residual need not be positive.
\par\noindent\emph{(E5): managed-return spectrum.}
Let
\[
T_K=\E[K(\cdot,Z)\otimes K(\cdot,Z)],
\quad\Sigma_{\Hh_K}=\E[X_t\otimes X_t].
\]
Because $\Sigma_t^0\prec\kappa_0I_N$ and
$\|h_t\|^2\leq\kappa_0$, we have
\[
\sigma_\varepsilon^2I_N\preceq\Sigma_t\preceq2\kappa_0I_N.
\]
For every $V\in\Hh_K$, conditional integration and the reproducing
property yield
\begin{align*}
\langle V,\Sigma_{\Hh_K}V\rangle_{\Hh_K}
&=\E\left[\frac1{N^2}\sum_{i,j}
V(Z_{i,t})\Sigma_{ij,t}V(Z_{j,t})\right],\\
\frac{\sigma_\varepsilon^2}{N}\E[V(Z)^2]
&\leq\langle V,\Sigma_{\Hh_K}V\rangle_{\Hh_K}
\leq\frac{2\kappa_0}{N}\E[V(Z)^2].
\end{align*}
Thus
\[
\frac{\sigma_\varepsilon^2}{N}T_K
\preceq\Sigma_{\Hh_K}
\preceq\frac{2\kappa_0}{N}T_K.
\]
The uniform characteristic law on the three-dimensional cube and the
Mat\'ern-$3/2$ equivalence $\Hh_K\simeq H^3(\mathcal Z)$ imply that
the positive eigenvalues $\tau_j$ of $T_K$ obey
\[
c_Kj^{-2}\leq\tau_j\leq C_Kj^{-2},\qquad j\geq1,
\]
by the Sobolev embedding spectral asymptotics (embedding singular
values have order $j^{-3/3}$). The min--max principle therefore gives
\[
\frac{\sigma_\varepsilon^2c_K}{N}j^{-2}
\leq\mu_j(\Sigma_{\Hh_K})
\leq\frac{2\kappa_0C_K}{N}j^{-2},\qquad j\geq1.
\]
In particular, $b=2>1$. These are bounds on the managed-return
second-moment operator, not on the centered return covariance or just
on the kernel itself.
\par\noindent\emph{(E6): positive and bounded optimal Sharpe.}
By Proposition~\ref{prop:mc_sharpe},
\[
0<p^\star=\frac{N}{\kappa_0}(v_\alpha+v_\beta L)<1,
\qquad\SR^\star=\sqrt{\frac{p^\star}{1-p^\star}}\in(0,\infty).
\]
Neither $p^\star$ nor $\SR^\star$ depends on $\rho$, so common
strictly positive finite Sharpe bounds hold over the entire stated
family. Under the calibration $p^\star=3/7$ the annualized marginal
Sharpe is exactly three for every $\rho$.
Finally, the pricing residual $m^\star_{t+1}=1-R^p_{t+1}(W^\star)$
satisfies
\[
|m^\star_{t+1}|\leq1+RB_X=:M_0<\infty,
\qquad\E[m^\star_{t+1}X_t\mid\mathscr F_t]=0.
\]
Thus it satisfies bounded conditional Bernstein noise inequalities;
for example,
\[
\E\left[\exp\left(\frac{|m^\star_{t+1}|}{M_0}\right)
-\frac{|m^\star_{t+1}|}{M_0}-1\ \middle|\mathscr F_t\right]
\leq e-2<1.
\]
All conditions are therefore verified for the specified family.
\end{proof}



\begin{proposition}[Functional Orthogonality and Exact Portfolio Regret]
\label{prop:portfolio_regret}

Let $\mathcal H$ be a linear space of admissible portfolio
policies, and define
\[
W_t(f)
=
\frac{1}{N}
\bigl(f(Z_{1,t}),\ldots,f(Z_{N,t})\bigr)'.
\]

Consider the population criterion
\[
Q(f)
=
\mathbb E\left[
\left(1-W_t(f)'R^e_{t+1}\right)^2
\right],
\]
and suppose that its minimum over $\mathcal H$ is attained at
$f^\star$.

Define the conditional moments
\[
\mu_i(Z_t)
=
\mathbb E[R^e_{i,t+1}\mid Z_t],
\qquad
\Omega_{ij}(Z_t)
=
\mathbb E[R^e_{i,t+1}R^e_{j,t+1}\mid Z_t].
\]

Assume that the relevant second moments are finite.

Then the population-optimal policy satisfies the
functional orthogonality condition
\begin{equation}
\boxed{
\sum_{i=1}^{N}
\mathbb E\left[
h(Z_{i,t})
\left(
\mu_i(Z_t)
-
\frac{1}{N}
\sum_{j=1}^{N}
\Omega_{ij}(Z_t)f^\star(Z_{j,t})
\right)
\right]
=0
}
\label{eq:functional_orthogonality}
\end{equation}
for every $h\in\mathcal H$.

Consequently, for every $f\in\mathcal H$,
\begin{equation}
\boxed{
\begin{aligned}
Q(f)-Q(f^\star)
=
\frac{1}{N^2}
\sum_{i,j=1}^{N}
\mathbb E\Big[
&\bigl(f(Z_{i,t})-f^\star(Z_{i,t})\bigr)\\
&\times
\bigl(f(Z_{j,t})-f^\star(Z_{j,t})\bigr)
\Omega_{ij}(Z_t)
\Big].
\end{aligned}
}
\label{eq:exact_portfolio_regret}
\end{equation}

In particular, this identity applies to the estimated policy
$\widehat f_{\lambda,T}$, conditionally on its training sample,
when population risk is evaluated on a fresh observation.
\end{proposition}


\begin{proof}

By conditioning on $Z_t$ and expanding the square,
the population criterion can be written as
\begin{equation}
\begin{aligned}
Q(f)
={}&1
-\frac{2}{N}
\sum_{i=1}^{N}
\mathbb E\left[
f(Z_{i,t})\mu_i(Z_t)
\right]\\
&+
\frac{1}{N^2}
\sum_{i,j=1}^{N}
\mathbb E\left[
f(Z_{i,t})f(Z_{j,t})
\Omega_{ij}(Z_t)
\right].
\end{aligned}
\label{eq:population_loss_expansion}
\end{equation}

Consider an arbitrary direction $h\in\mathcal H$.
Since $\mathcal H$ is linear,
$f^\star+\varepsilon h\in\mathcal H$ for every
$\varepsilon\in\mathbb R$.

The Gâteaux derivative of $Q$ at $f^\star$
in direction $h$ is
\[
DQ(f^\star)[h]
=
\left.
\frac{d}{d\varepsilon}
Q(f^\star+\varepsilon h)
\right|_{\varepsilon=0}.
\]

Differentiating \eqref{eq:population_loss_expansion}
and using the symmetry
$\Omega_{ij}(Z_t)=\Omega_{ji}(Z_t)$ gives
\[
\begin{aligned}
DQ(f^\star)[h]
={}&
-\frac{2}{N}
\sum_{i=1}^{N}
\mathbb E[
h(Z_{i,t})\mu_i(Z_t)
]\\
&+
\frac{2}{N^2}
\sum_{i,j=1}^{N}
\mathbb E[
h(Z_{i,t})f^\star(Z_{j,t})
\Omega_{ij}(Z_t)
].
\end{aligned}
\]

Since $f^\star$ minimizes $Q$ over the linear space
$\mathcal H$, we have
\[
DQ(f^\star)[h]=0
\qquad
\text{for every }h\in\mathcal H.
\]

Rearranging establishes
\eqref{eq:functional_orthogonality}.

We now establish the exact regret identity.
For any $f\in\mathcal H$, expanding the difference
between the two population losses yields
\[
\begin{aligned}
Q(f)-Q(f^\star)
={}&
-\frac{2}{N}
\sum_i
\mathbb E[
(f(Z_{i,t})-f^\star(Z_{i,t}))
\mu_i(Z_t)
]\\
&+
\frac{1}{N^2}
\sum_{i,j}
\mathbb E[
(f(Z_{i,t})f(Z_{j,t})
-f^\star(Z_{i,t})f^\star(Z_{j,t}))
\Omega_{ij}(Z_t)
].
\end{aligned}
\]

Using the algebraic identity
\[
\begin{aligned}
f(Z_{i,t})f(Z_{j,t})
-f^\star(Z_{i,t})f^\star(Z_{j,t})
={}&
(f(Z_{i,t})-f^\star(Z_{i,t}))
(f(Z_{j,t})-f^\star(Z_{j,t}))\\
&+
(f(Z_{i,t})-f^\star(Z_{i,t}))
f^\star(Z_{j,t})\\
&+
f^\star(Z_{i,t})
(f(Z_{j,t})-f^\star(Z_{j,t})),
\end{aligned}
\]
and exploiting the symmetry of $\Omega$, we obtain
\[
\begin{aligned}
Q(f)-Q(f^\star)
={}&
\frac{1}{N^2}
\sum_{i,j}
\mathbb E\Big[
(f(Z_{i,t})-f^\star(Z_{i,t}))
(f(Z_{j,t})-f^\star(Z_{j,t}))
\Omega_{ij}(Z_t)
\Big]\\
&-
\frac{2}{N}
\sum_i
\mathbb E\left[
(f(Z_{i,t})-f^\star(Z_{i,t}))
\left(
\mu_i(Z_t)
-
\frac{1}{N}
\sum_j
\Omega_{ij}(Z_t)f^\star(Z_{j,t})
\right)
\right].
\end{aligned}
\]

Because $\mathcal H$ is linear,
$f-f^\star\in\mathcal H$.
Taking $h=f-f^\star$ in
\eqref{eq:functional_orthogonality}
shows that the second term vanishes.

Therefore,
\[
\boxed{
Q(f)-Q(f^\star)
=
\frac{1}{N^2}
\sum_{i,j=1}^{N}
\mathbb E\left[
(f(Z_{i,t})-f^\star(Z_{i,t}))
(f(Z_{j,t})-f^\star(Z_{j,t}))
\Omega_{ij}(Z_t)
\right].
}
\]

This proves the result.
\end{proof}




\begin{equation}
\label{eq:empirical_portfolio_regret}
\boxed{
\begin{aligned}
\widehat{\mathcal R}_T
=\frac{1}{S}\sum_{s=1}^{S}\Bigg[
&\left(1-\widehat W_T^{(s)\prime}R^{(s)}\right)^2\\
-&\left(1-W^{\star\prime}(Z^{(s)})R^{(s)}\right)^2
\Bigg].
\end{aligned}
}
\end{equation}
